\chapter{Patterns: optimistic updates}
\label{ch:optimistic}

\begin{aside}
Wait for the server before updating the UI? Too slow. Update
first, reconcile later. The user feels fast; the network does
its work in the background.
\end{aside}

\section{The pattern}

\begin{enumerate}[leftmargin=*]
  \item User clicks like button.
  \item UI immediately shows liked.
  \item Background request posts to server.
  \item On success: confirm.
  \item On failure: revert with toast.
\end{enumerate}

\section{Tracking pending state}

Each optimistic update carries a state: \code{pending},
\code{confirmed}, \code{reverted}.

\begin{lstlisting}
struct OptimisticUpdate {
    Update u;
    enum State { Pending, Confirmed, Reverted } state;
};
\end{lstlisting}

\section{Visible pending}

A subtle visual cue: dim the optimistic value, or show a
spinner. Or nothing --- if reverts are rare.

\section{Conflict with other updates}

If two optimistic updates affect the same value, the order
matters. Apply in order; revert in reverse on failure.

\section{Idempotency}

Server-side: idempotent updates make retries safe. A
client that thinks it failed but actually succeeded can
retry without duplicating.

\section{Recap}

\begin{recap}
\begin{itemize}[leftmargin=*]
  \item Update local first; reconcile after.
  \item Track pending state.
  \item Subtle visual cue, optional.
  \item Idempotent server-side for safe retries.
\end{itemize}
\end{recap}

\section*{Workshop}

\begin{enumerate}[leftmargin=*]
  \item Implement an optimistic like button.
  \item Add revert-on-failure with toast.
  \item Test with network throttling.
\end{enumerate}
